\section{Learning C-Trees}\label{sec:v8-ctree-math}

A C-Tree has two learnable operators: a state map \(g\) that
summarizes a span, and a scorer \(f\) that reads a state and
returns a number. Neural operators are a natural representation
family for both: a transformer is a specific class of neural
operator. We instantiate \(f\) and \(g\) as prompted LLMs. We use a prompted LLM rather than another neural
operator because the produced state is human-readable text, which
makes node-level audit possible: an analyst (or a sampled expert)
can inspect a leaf summary, a re-summarized state, or a merged
parent state and decide whether the local laws hold. The synthetic
Markov sandbox in Section~\ref{sec:v10-controlled-checks} keeps the
operators simple to make the budget question visible; the
Manifesto LLM replication uses the same construction with prompted
LLMs.

This section defines the C-Tree pieces, states the three local
laws \(g\) must satisfy, and proves what those laws buy at the
root. Throughout, the document and the realized tree are fixed.

\subsection{Objects and States}

Let \(\mathcal X\) be the input domain of documents and document
spans, let \(\mathcal Y\) be the task-output space, and let
\(\mathcal Z\) be the state space. In an LLM instantiation,
\(\mathcal X\) and \(\mathcal Z\) are both spaces of text:
\(\mathcal X\) is raw document text and \(\mathcal Z\) is short
evidence summaries written by \(g\). For the Manifesto LLM
replication, \(\mathcal X\) is party-platform text and \(\mathcal Y\)
is a seven-point policy score on one dimension (or the vector of six
dimension scores when one tree serves all six).

Fix a document \(x\in\mathcal X\). A realized C-Tree is a finite
ordered binary tree \(T\) with node set \(V(T)\) and leaves forming a
fixed contiguous partition of \(x\). Each node \(v\in V(T)\) represents
a span \(x_v\in\mathcal X\): at a leaf, \(x_v\) is the raw document
span; at an internal node with left child \(u\) and right child \(w\),
\(x_v=x_u\concat x_w\). Here \(\concat\) denotes ordered concatenation.
The root span is \(x_{\mathrm{root}}=x\).

The C-Tree learns one state map, written \(g\). At a leaf, \(g\) reads
a raw document span and returns an evidence state. At an internal
node, \(g\) reads the two child states and returns the parent state.
The same \(g\) can also be applied to a state that \(g\) already
produced, e.g.\ during a cleanup pass that rewrites a state before
it merges with its sibling. The stored state at node \(v\) is
\[
  z_v =
  \begin{cases}
    g(x_v), & v \text{ is a leaf},\\
    g(z_u,z_w), & v \text{ has children }u,w.
  \end{cases}
\]

\begin{figure}[t]
\centering
\includegraphics[width=0.85\linewidth]{01_base_plain.pdf}
\caption{\textbf{C-Tree summary tree.} Raw blocks
\(b_1,\ldots,b_8\) at the bottom; \(s_i = g(b_i)\) is the leaf
summary; \(s_{ij}=g(s_i,s_j)\) and so on up to the root summary
\(s_{\mathrm{root}}\). One learned operator \(g\) handles every
arrow: leaf encoding, re-summarization, and merge.}
\label{fig:v8-base-tree}
\end{figure}

Figure~\ref{fig:v8-base-tree} shows the resulting summary tree on
an eight-block document. The same operator \(g\) handles every
arrow; only the input shape changes (one raw block at a leaf, one
already-produced state for re-summarization, two states at a
merge).

A scorer reads one state and returns a number. We write \(f(z)\) for
a scorer applied to state \(z\), \(f^\ast\) for the target oracle
(the expert or rubric scorer the analyst wants to recover), and
\(\widehat f\) for the fitted cheaper proxy used for dense local
supervision. In an LLM instantiation, the scorer is a prompt that
places the state where it expects the document, and the proxy
\(\widehat f\) is the same prompt applied to a smaller LLM.
\(f^\ast\) is observed only when the oracle is called.

For any scorer \(q\), query equivalence is
\[
  a\sim_q b
  \quad\Longleftrightarrow\quad
  d_Y(q(a),q(b))=0 .
\]
Here \(d_Y\) is the task-space discrepancy on \(\mathcal Y\), and
\(a\sim_q b\) means that \(a\) and \(b\) give the same answer under
\(q\). We specialize to the oracle query \(q=f^\ast\) and write
\[
  a\sim_\ast b
  \quad\Longleftrightarrow\quad
  d_Y(f^\ast(a),f^\ast(b))=0 .
\]
This convention treats an \(f^\ast\) call as observation-error-free
once it is made; the estimation problem is that such calls are sparse.
Section~\ref{sec:v8-audit-certificate} replaces the equality with a
small tolerance and bounds the resulting drift in the reported number.

\begin{table}[t]
  \centering
  \small
  \begin{tabular}{@{}p{0.22\linewidth}p{0.68\linewidth}@{}}
    \toprule
    Symbol & Meaning \\
    \midrule
    \(x,T,V(T),v\) & Document, realized C-Tree, node set of the tree, and a node. \\
    \(u,w\) & Left and right child nodes of an internal node. \\
    \(x_v,z_v\) & Document span represented by node \(v\), and the state stored at \(v\). \\
    \(g\) & One learned C-Tree state map/program, called on raw leaves, on already-produced states, and on pairs of child states. \\
    \(\mathcal X,\mathcal Y,\mathcal Z,\concat\) & Input domain, task-output space, state space, and ordered concatenation. \\
    \(d_Y,\sim_\ast\) & Task-space discrepancy and oracle equivalence. \\
    \(f^\ast,\widehat f\) & Target oracle readout and fitted proxy readout. \\
    C1/C2/C3 & Sufficiency, idempotence, and merge consistency (the local laws). \\
    \bottomrule
  \end{tabular}
  \caption{\textbf{Core notation for the C-Tree math.}}
  \label{tab:v8-core-notation}
\end{table}

\begin{assumption}[Fixed Realized Tree]\label{ass:v8-fixed-tree}
For each document \(x\), the leaf partition and tree topology are conditioned
on as fixed when stating the local laws and the objective.
\end{assumption}

The assumption isolates representation validity from partition
selection. A deterministic partition rule returns the same statement
after conditioning on its realized tree. The leaf partition is a
chunker design choice; in the Manifesto LLM replication it is a
fixed token budget per leaf, swept from \(256\) to \(8{,}192\) tokens
(Section~\ref{sec:v8-benoit-tree-setup}).

\subsection{Local Laws}

The three local laws are stated on the realized calls made by the tree.

\begin{definition}[Oracle C-Tree Local Laws]\label{def:v8-local-laws}
For a realized tree \(T\), state map \(g\), and target oracle \(f^\ast\), the
local laws are:
\[
\begin{aligned}
\mathrm{C1:}\quad&
g(x_v)\sim_\ast x_v
&&\text{for every leaf }v,\\
\mathrm{C2:}\quad&
g(z_v)\sim_\ast z_v
&&\text{for every produced state }z_v,\\
\mathrm{C3:}\quad&
g(z_u,z_w)\sim_\ast x_u\concat x_w
&&\text{for every internal node }v=(u,w).
\end{aligned}
\]
\end{definition}

\begin{figure}[t]
\centering
\includegraphics[width=0.95\linewidth]{07_local_laws_plain.pdf}
\caption{\textbf{The three local laws on a realized tree.}
C1 (sufficiency) is checked at a leaf:
\(f^\ast(b_3)\approx f^\ast(s_3)\)? C2 (idempotence) is checked at
any produced state:
\(f^\ast(s_{1234})\approx f^\ast(g(s_{1234}))\)? C3 (merge) is
checked at any internal node:
\(f^\ast(s_{56})\approx f^\ast(b_5\!\oplus\!b_6)\)? The same
operator \(g\) is the subject of all three checks; the laws ask
whether each kind of call lands in the same oracle fiber as the
span it represents.}
\label{fig:v8-local-laws-tree}
\end{figure}

C1 (sufficiency) says that once a raw document span enters the tree,
the state \(g(x_v)\) is already sufficient for oracle-facing
questions about that span. A leaf state may be shorter, more
structured, or more readable than the raw text, but it must stay in
the same \(f^\ast\)-fiber.

C2 (idempotence) says that summarizing a state again gives the same
oracle-facing answer as the state itself: a summary of a summary
lands in the same \(f^\ast\)-fiber as the summary. This prevents a
cleanup pass from turning a conditional claim into a different
oracle-facing position.

C3 (merge consistency) says that merging two child states produces
a parent state equivalent to the represented parent span
\(x_u\concat x_w\). The parent state must preserve interactions
that are invisible from scalar child scores: qualifications,
contradictions, cross-references, and tradeoffs that span across
two child blocks.

The local laws define which learned state maps count as sensible
for the task; they double as post hoc diagnostics. The exact
statement above assumes \(f^\ast\) can be called at every node,
which is too expensive in deployment. So dense local training uses
the proxy \(\widehat f\) at every node, and a small set of sampled
\(f^\ast\) calls then estimates and corrects the gap between the
proxy and the oracle. Section~\ref{sec:v8-objective} writes the
objective for this setup; Section~\ref{sec:v8-audit-certificate}
writes the correction.

\subsection{Why States Matter}

The local laws are state laws, not scalar-score laws. If a leaf
stores only \(f^\ast(x_v)\), then the parent has no access to
evidence that may matter only after two spans are read together.
Two child spans can have the same local score and still differ in
whether they supply a qualification, condition, or exception that
the parent score depends on. The state \(z_v\) is the carrier for
that interaction information.

A C-Tree can support more than one downstream readout of the same
compressed document, provided those readouts respect the oracle
equivalence relation. The theorem in the next subsection states
that if the states satisfy the local laws, then the root state
lies in the same \(\sim_\ast\) class as the original document.
Anything that is constant on those classes can then be evaluated
at the root.

\subsection{Root Preservation}

\begin{theorem}[Local Laws Preserve the Root]\label{thm:v8-root-preservation}
If C1, C2, and C3 hold exactly for \(g\), \(T\), and \(f^\ast\), then the root
state preserves the document oracle:
\[
  d_Y\!\left(f^\ast(z_{\mathrm{root}}),f^\ast(x)\right)=0 .
\]
For stochastic C-Tree reductions and \(m\ge 1\) rounds of
re-summarizing produced states, the expected oracle distortion is
zero.
\end{theorem}

\begin{proof}[Proof sketch]
Assume C1, C2, and C3 hold exactly on the realized tree \(T\) for
\(g\) and \(f^\ast\). The argument is induction on the depth of a
node. At a leaf, C1 says the leaf state \(g(x_v)\) gives the same
oracle answer as the raw leaf span \(x_v\). At an internal node with
children whose stored states already give the same oracle answer as
their represented spans, C3 says the merge \(g(z_u,z_w)\) gives the
same oracle answer as the concatenated parent span \(x_u\concat x_w\).
The single bookkeeping case is when a state is summarized again
before it merges; C2 says re-summarizing does not change the oracle
answer either, so the induction continues through that step. The
root case is the same identity at the root node:
\(f^\ast(z_{\mathrm{root}})=f^\ast(x)\). For stochastic reductions,
the same identity holds for every realization, so it survives the
expectation. The full induction is in
Appendix~\ref{app:v8-full-proofs}.
\end{proof}

The theorem says the root state is interchangeable with the raw
document for any quantity that depends only on the oracle answer.
That interchangeability is the load-bearing claim for everything
downstream: training, scoring, and pairwise comparison can use the
root state in place of the original document.

\subsection{Pairwise and Cross-Document Readouts}

\begin{corollary}[Oracle-Compatible Readouts Pass Through C-Tree Nodes]
\label{cor:v8-preferences-pass-to-root}
Call a readout \(\Phi\) \emph{oracle-compatible} if
\[
  a\sim_\ast b \quad\Longrightarrow\quad \Phi(a)=\Phi(b).
\]
Under exact C1/C2/C3 local laws, every realized node \(v\in V(T)\) satisfies
\[
  \Phi(z_v)=\Phi(x_v).
\]
Taking \(v=\mathrm{root}\) gives the root statement
\(\Phi(z_{\mathrm{root}})=\Phi(x)\).

For pairwise preferences, let \(P\) be a comparison rule that is
oracle-compatible in both arguments:
\[
  a\sim_\ast a',\ b\sim_\ast b'
  \quad\Longrightarrow\quad
  P(a,b)=P(a',b').
\]
Then any two valid node states, from the same C-Tree or from different
C-Trees satisfying the same oracle laws, can replace their raw blocks:
\[
  P(z_v,z_{v'})=P(x_v,x_{v'}).
\]
\end{corollary}

\begin{proof}[Proof sketch]
Assume C1, C2, and C3 hold exactly as in
Theorem~\ref{thm:v8-root-preservation}, plus that the readout
\(\Phi\) is oracle-compatible. Pick any node \(v\) and read the
subtree rooted at \(v\) as a smaller C-Tree on the span \(x_v\).
Theorem~\ref{thm:v8-root-preservation} applied to that subtree gives
the same oracle answer at the subtree's root, which is the stored
state \(z_v\): in symbols, \(z_v\sim_\ast x_v\). Oracle-compatibility
of \(\Phi\) turns this equivalence into the equality
\(\Phi(z_v)=\Phi(x_v)\). The pairwise version is the same step
applied to each argument of \(P\) separately, because \(P\) is
assumed oracle-compatible in both arguments. The full argument is in
Appendix~\ref{app:v8-full-proofs}.
\end{proof}

The node-level form makes the local-law losses useful as local
supervision, beyond serving as an indirect route to a final document
score. If the raw block \(x_v\) is a meaningful unit for a
task-facing readout, then a valid state \(z_v\) is meaningful for
the same readout, provided the readout depends only on the oracle
equivalence class. The root is the largest such block. The same
applies across documents: two states from different documents can
replace their raw spans for any comparison that depends on the
oracle, as long as both trees satisfy the oracle laws on the same
target.

A preference about information not represented in \(f^\ast\) is not
certified by this theorem, even if the root preserves \(f^\ast\)
exactly. To certify that preference, the oracle must be expanded to
include the relevant readout, or the local laws must be stated for a
richer query family. Within the chosen oracle interface, the
corollary is the preference-learning consequence of the exact
theorem: a valid C-Tree root can stand in for the original long
document for scores that respect that interface.

The state story connects to classical mergeable summaries
(Section~\ref{sec:v8-scope-related} develops the connection). A
classical sketch chooses the state by hand and proves the merge law.
A learned C-Tree chooses a state class and uses local-law losses
plus audits to decide whether the fitted state behaves as a valid
task-relative state machine.

\subsection{From Exact Laws to Learning}

Theorem~\ref{thm:v8-root-preservation} describes the endpoint: if
every realized leaf, state re-entry, and merge call stays in the
right oracle fiber, the root state can stand in for the full
document. Learning starts from the fact that this endpoint is
only partially observable. If \(f^\ast\) could be called at every leaf, state,
and merge, the learning problem would be almost direct: search
for a \(g\) whose C1/C2/C3 discrepancies are small under
\(f^\ast\), while also fitting the document-level expert target.

In a deployed setting, oracle calls are not free at every node. An
\(f^\ast\) call may be an expert label, an expensive teacher
judgment, or an aggregate label derived from application-aligned
coding. The oracle local-law loss is therefore observable at only a
small set of sampled nodes. C-TreePO optimizes one learned state
map \(g\) for two purposes at once: it should be useful at the
root, and it should behave like a valid C-Tree state throughout the
realized tree.

A realized tree makes this approximation problem concrete. At a
leaf, we ask whether the state preserves the raw evidence. At a
produced state, we ask whether summarizing it again changes the
oracle-facing content. At an internal node, we ask whether merging
child states preserves the concatenated span. We write the loss
that pushes \(g\) toward those three answers next.
